The perturbation theorem for Lorentz gases implies continuity in $R$ of the
spectral projectors and leading eigenvalues.  Hence, on compact preparation
sets, the maps
\[
 (R,q)\mapsto j_{R,q},\quad \bar\tau_{R,q},\quad v_{R,q},\quad D_{R,q}
\]
are continuous.  At $q=0$, Paper I gives the explicit nonzero derivative of
$\bar\tau_R^{\rm L}$.  No differentiability of the full tilted spectrum with
respect to $R$ is used in this paper.

\section{Main theorem and scope}\label{sec:main}
\begin{theorem}[Same-platform physical transport theorem]\label{thm:main}
The natural Liouville current preparation, actual cell displacement, actual
free-flight time, ballistic velocity, comoving Brownian fluctuation, physical
diffusion coefficient, and exponential random-horizon control all arise from
the single periodic Sinai family of Paper I.  There is no microscopic random
forcing and no surrogate clock.
\end{theorem}
\begin{proof}
Combine \cref{thm:root,thm:asip,thm:physical-limit,cor:sde,thm:window}.
\end{proof}

\section*{Truth boundary}
The physical diffusion describes fluctuations about the current-prepared
ballistic velocity.  The theorem is local in the canonical field and uses the
standard Lorentz spectral/ASIP packet.  External review of the joint
pressure-root and uniform random-window arguments remains pending.
\printbibliography
